\documentclass{beamer}

\usepackage{pgfplots}
\usepackage{pgf-pie}
\usepackage{graphicx}

\usepgfplotslibrary{statistics}
\pgfplotsset{compat=1.18}

\definecolor{dark_gray}{HTML}{121111}
\definecolor{bright_orange}{HTML}{ffb717}

\definecolor{pastel_yellow}{HTML}{fccb6f}
\definecolor{pastel_blue}{HTML}{79c6fc}
\definecolor{verde_lima}{HTML}{77fc65}
\definecolor{rosado_gay}{HTML}{e05a89}
\definecolor{cafe_claro}{HTML}{e0905a}
\definecolor{purpl_pikmin}{HTML}{a24ff0}

\setbeamercolor{background canvas}{bg=dark_gray}
\setbeamercolor{normal text}{fg=white}
\setbeamercolor{title}{fg=bright_orange}
\setbeamercolor{frametitle}{fg=bright_orange}
\setbeamercolor{structure}{fg=bright_orange}

\setbeamertemplate{frametitle}{

    %\hspace*{-0.8cm}
    \centering
    \insertframetitle
    \vspace{0.2cm}

    \nointerlineskip
    \hspace*{-2cm}
    {\color{bright_orange}\rule{\dimexpr\paperwidth+1cm}{2pt}}

    \vspace{0.1cm}
}

\setbeamertemplate{footline}{
    \begin{beamercolorbox}[wd=\paperwidth, ht=2pt]{footline}
        \color{bright_orange}\rule{\paperwidth}{2pt}
    \end{beamercolorbox}
    \vspace{0.4cm}
}

\newcommand{\StyleError}{\color{red}}
\newcommand{\StyleKeyword}{\color{pastel_blue}}
\newcommand{\StyleIdentifier}{\color{white}}
\newcommand{\StyleLiteral}{\color{pastel_yellow}}
\newcommand{\StyleString}{\color{verde_lima}}
\newcommand{\StyleOperator}{\color{rosado_gay}}
\newcommand{\StylePunctuator}{\color{cafe_claro}}

% para usar los estilos
\newcommand{\Error}[1]{{\StyleError #1}}
\newcommand{\Keyword}[1]{{\StyleKeyword #1}}
\newcommand{\Identifier}[1]{{\StyleIdentifier #1}}
\newcommand{\Literal}[1]{{\StyleLiteral #1}}
\newcommand{\String}[1]{{\StyleString #1}}
\newcommand{\Operator}[1]{{\StyleOperator #1}}
\newcommand{\Punctuator}[1]{{\StylePunctuator #1}}

\setbeamertemplate{navigation symbols}{}

% MAXIMO LINEAS POR SLIDE
\newcommand{\SourceLinesPerSlide}{20}

\newcommand{\SourceFrameTitle}{Source code}
\newcommand{\SourceFont}{\ttfamily\scriptsize} % font
\newcommand{\SrcTab}{\ \ \ \ } % tabs
\newcommand{\SrcWrapIndent}{2em} % sangría de las líneas que continúan
\newcommand{\SrcBreak}{\penalty 20 } % lo mismo

\newcounter{srcvisual} % líneas en slide
\newcounter{srcpart} % numero de diapositiva
\newtoks\srcbuffer % buffer para linea
\newsavebox{\srcbox} % cada para medir tamaño del codigo


% Configuración para mostrar las líneas de código.
\newcommand{\SrcSetup}{
    \SourceFont \hyphenchar\font=-1
    \hsize=\textwidth \parindent=0pt
    \tolerance=10000 \pretolerance=-1
    \rightskip=0pt plus 1fil
    \hangindent=\SrcWrapIndent \hangafter=1
}

%mostrar linea de codigo
\newcommand{\SrcRenderLine}[1]{{\SrcSetup \strut #1\par}}

% medir cuanto espacio ocupa una linea
\newcommand{\SrcMeasure}[1]{%
    \setbox\srcbox\vbox{\SrcRenderLine{#1}\xdef\SrcNL{\the\prevgraf}}}

\newcommand{\SrcEmitFrame}{%
    \stepcounter{srcpart}%
    \begin{frame}[t]{\SourceFrameTitle{} (\arabic{srcpart})}
        \the\srcbuffer
    \end{frame}%
    \global\srcbuffer{}%
    \setcounter{srcvisual}{0}%
}

% crea slide con lo q hay en buffer
\newcommand{\SrcFlush}{\ifnum\value{srcvisual}>0\expandafter\SrcEmitFrame\fi}

% para cada linea de codigo en source
\newcommand{\SrcLine}[1]{%
    \SrcMeasure{#1}%

    % si no cabe linea completa
    \ifnum\numexpr\value{srcvisual}+\SrcNL\relax>\SourceLinesPerSlide\relax
        \expandafter\SrcFlush
    \fi
    \global\srcbuffer\expandafter{\the\srcbuffer\SrcRenderLine{#1}}%
    \addtocounter{srcvisual}{\SrcNL}%
    \ifnum\value{srcvisual}<\SourceLinesPerSlide\relax\else\expandafter\SrcFlush\fi
}

% comando d input para el codigo fuente
\newcommand{\ShowSourceCode}[1]{%
    \setcounter{srcvisual}{0}\setcounter{srcpart}{0}\global\srcbuffer{}%
    \input{#1}%
    \SrcFlush
}

\begin{document}

\title{Project 01}
\subtitle{Compilers \& Interpreters}
\author{Javier Lee, Steven Loaíciga, Luan Valverde}
\date{2026, semester II}

\setbeamertemplate{headline}{
    \vspace{0.4cm}
    \begin{beamercolorbox}[wd=\paperwidth, ht=2pt]{footline}
        \color{bright_orange}\rule{\paperwidth}{2pt}
    \end{beamercolorbox}
}

\begin{frame}
    \maketitle
\end{frame}

\setbeamertemplate{headline}{}

\begin{frame}{Scanning process}
    Scanning, or \textit{lexical analysis}, is \textit{conceptually} the first phase in the compilation process.
    \begin{itemize}
        \item The \textit{scanner} receives the \textit{source code} as input.
        \item It reads the entire file, character by character.
        \item It groups these characters into \textit{tokens}.
        \item These tokens are then sent to the parser via the getToken() function.
    \end{itemize}
    \centering
    \includegraphics[width=0.8\textwidth]{figures/scannerimg.png}
% add image !!!!

\end{frame}

\begin{frame}{getToken() function}
    This function is the main ----- of the scanner.
    \begin{itemize}
        \item The entire program is a string stored in the computer's memory.
        \item The scanner keeps track of the current position.
        \item It scans through the string until it forms an entire token.
        \item It skips what is considered as "spaces" by the specific language...
            \begin{itemize}
                \item Line jumps, tabs, comments, ...
            \end{itemize}
    \end{itemize}
    \centering
\end{frame}

\begin{frame}{Tokens}
    Tokens are \textit{usually} a data structure with two fields or attributes: 
    \begin{itemize}
        \item It's \textit{token code}, which is a number that specifies if it is a variable, constant, comma, etc.
        \item The \textit{lexeme}, which is usually a pointer to the particular string it refers to.
    \end{itemize}
    \centering
\end{frame}

\begin{frame}{Flex}
    Flex is a toll for generating \textit{lexical analyzers} (scanners).
    \begin{itemize}
        \item It lets the user define the patterns that represent tokens using \textit{regular expressions}.
        \item So instead of writing the scanner from scratch, Flex generates the corresponding C code.
    \end{itemize}
    \centering
\end{frame}

\begin{frame}{Flex (cont.)}
    \begin{itemize}
        \item Flex receives a file with the \textit{.l} extension.
        \item This file contains the patterns and corresponding actions of the language.
        \item Flex processes this file and generates a C file called \textit{lex.yy.c}
        \item This file is the scanner, and it can be modified further by the user if they deem it necessary.
    \end{itemize}
\end{frame}

% Múltiples slides con el programa fuente que le entró a la fase de scanning (i.e., después del preproceso)

\begin{frame}{Source code}
    The following slides contain the source code that was processsed by the scanner.
\end{frame}

\ShowSourceCode{latex_aux/source_code.tex}

% ...

% Histograma de las cantidades de cada tipo de token encontrados en el fuente.

\begin{frame}{Statistics - Histogram}
\centering
\begin{tikzpicture}

\begin{axis}[
    ybar,
    width=0.85\textwidth,
    height=0.55\textheight,
    bar width=30pt,
    ylabel={Quantity},
    symbolic x coords={
        Keyword,
        Identifier,
        Literal,
        String,
        Operator,
        Punctuator
    },
    xtick=data,
    xtick style={draw=none},
    ytick style={draw=white},
    axis lines=left,
    axis line style={-},
    ymin=0,
    enlarge y limits={upper,value=0.30},
    enlarge x limits=0.20,
    nodes near coords,
    xticklabel style={
        rotate=30,
        anchor=east
    },
    tick label style={font=\small},
]

\addplot[ fill=pastel_yellow, draw=pastel_yellow
] table[
    x=Category,
    y=Quantity,
    col sep=space
] {latex_aux/token_stats.dat};


\end{axis}

\end{tikzpicture}

\end{frame}

% Gr´afico de pastel de las categor´ıas l´exicas encontradas (misma informaci´on que el histograma).

\begin{frame}{Statistics - Pie chart}
\begin{center}
\input{latex_aux/pie_chart.tex}
\end{center}
\end{frame}

% finall

\setbeamertemplate{headline}{
    \vspace{0.4cm}
    \begin{beamercolorbox}[wd=\paperwidth, ht=2pt]{footline}
        \color{bright_orange}\rule{\paperwidth}{2pt}
    \end{beamercolorbox}
}

\begin{frame}
    That's all. Thank you for reading this presentation.
\end{frame}

\end{document}